% !TeX TXS-program:compile = txs:///pdflatex

\documentclass{article}
\usepackage[margin=1cm,a4paper]{geometry}
\usepackage{customenvs-mathpictos}

\usepackage{tikz}
\usepackage{longtable}
\usepackage{xintexpr}
\usepackage{fancybox}
\usepackage{xcolor}
\usepackage{hyperref}
\pagecolor{teal!5}
\setlength\parindent{0pt}

\newcommand\tikzvalign[2]{%
	\begin{tikzpicture}
		\useasboundingbox (0,0) rectangle (1in,1in) ;
		\draw (0.5in,0.5in) node[text width=1in,align=center] {\texttt{#2}\\~\\{\tiny \texttt{(pages #1-\inteval{#1+4})}}} ; ;
	\end{tikzpicture}
}

\begin{document}

\pagestyle{empty}

\part*{Sub-package customenvs-mathpictos (0.1c)}

\section{French version}

Inspiré d'un travail de \textsf{Vincent Le Gruiec} (\url{https://www.vmaths.fr}).

\medskip

\textsf{[fr]} : \verb|\cemaththemeicon[includegraphics options]{bordures=...,fond=TF,couleur fond=...}{type}|

\def\tmpcptxint{1}

\begin{longtable}{|c|c|c|c|c|c|}
	\cline{2-6}
	\multicolumn{1}{c}{} & \multicolumn{5}{|c|}{\textbf{\texttt{bordures=}}} \\ \hline
	\textbf{\texttt{type}} & \texttt{aucune} {\footnotesize (défault)} & \texttt{11} & \texttt{10} & \texttt{01} & \texttt{00} \\ \hline
	%test génération
	\xintFor* #1 in {\xintCSVtoList{%
			binomiale,%
			convexite,%
			fctcube,%
			fctcarre,%
			variations,%
			probas,%
			nbderive,%
			ensembles,%
			lgn,%
			equadiff,%
			fcttrigo,%
			fctaffine,%
			integration,%
			intervalles,%
			limitesfct,%
			limitessuite,%
			calclitteral,%
			fctlog,%
			fctinv,%
			histogramme,%
			defsuites,%
			systemes,%
			chasles,%
			equations,%
			espace,%
			espaceproj,%
			pbgeom,%
			espacevect,%
			recurrence,%
			pourcentages,%
			fractions,%
			eqbarres,%
			scratch,%
			python,%
			tableur,%
			graphes,%
			arithm,%
			geogebra,%
			toilerecur,%
			cercletrigo,%
			pythagore,%
			thales,%
			outilsgeom,%
			numeration,%
			transfo,%
			applis,%
			booleen,%
			fctexpo,%
			pdtscal,%
			solide,%
			algo,%
			nwks,%
			gendarmes,%
			pixelart,%
			venn,%
			arbrebin,%
			matrices,%
			thmvalmoy,%
			pascal,%
			addposee,%
			arbrecalc,%
			sohcahtoa,%
			fibonnaci,%
			quatreops,%
			reglin,%
			moustaches,%
			fractale,%
			binaire,%
			patterns
		}}\do{%
		\tikzvalign{\tmpcptxint}{#1}&
		\raisebox{-\depth}{\cemaththemeicon[height=1in]{bordures=aucune}{#1}}&
		\raisebox{-\depth}{\cemaththemeicon[height=1in]{bordures=11}{#1}}&
		\raisebox{-\depth}{\cemaththemeicon[height=1in]{bordures=10}{#1}}&
		\raisebox{-\depth}{\cemaththemeicon[height=1in]{bordures=01}{#1}}&
		\raisebox{-\depth}{\cemaththemeicon[height=1in]{bordures=00}{#1}}\xdef\tmpcptxint{\inteval{\tmpcptxint+5}}\\ \hline
	}
\end{longtable}

\pagebreak

\section{English version}

Inspired by a worf of \textsf{Vincent Le Gruiec} (\url{https://www.vmaths.fr}).

\medskip

\textsf{[en]} : \verb|\cemaththemeicon[includegraphics options]{borders=...,bg=TF,colback=...}{type}|

\def\tmpcptxint{1}

\begin{longtable}{|c|c|c|c|c|c|}
	\cline{2-6}
	\multicolumn{1}{c}{} & \multicolumn{5}{|c|}{\textbf{\texttt{borders=}}} \\ \hline
	\textbf{\texttt{type}} & \texttt{none} {\footnotesize (default)} & \texttt{11} & \texttt{10} & \texttt{01} & \texttt{00} \\ \hline
	%test génération
	\xintFor* #1 in {\xintCSVtoList{%
			binom,%
			convexity,%
			cubicfct,%
			quadfct,%
			variations,%
			probas,%
			derivnb,%
			sets,%
			lln,%
			diffeq,%
			trigfct,%
			affinfct,%
			integration,%
			intervals,%
			fctlim,%
			seqlim,%
			littcalc,%
			logfct,%
			invfct,%
			histogram,%
			seqdef,%
			systems,%
			chasles,%
			equations,%
			space,%
			projspac,%
			geompb,%
			vectspace,%
			recurrence,%
			percentages,%
			fractions,%
			bareq,%
			scratch,%
			python,%
			spreadsheet,%
			graphs,%
			arithmetic,%
			geogebra,%
			cabweb,%
			trigcircle,%
			pythagora,%
			thales,%
			geomtools,%
			numeration,%
			transform,%
			applications,%
			boolean,%
			expfct,%
			dotproduct,%
			solid,%
			algorithm,%
			nwks,%
			squeezethm,%
			pixlart,%
			venn,%
			bintree,%
			matrix,%
			meanvaluethm,%
			pascal,%
			arithmop,%
			bincalctree,%
			sohcahtoa,%
			fibonnaci,%
			fourops,%
			linreg,%
			boxplot,%
			fractal,%
			binary,%
			patterns
	}}\do{%
		\tikzvalign{\tmpcptxint}{#1}&
		\raisebox{-\depth}{\cemaththemeicon[height=1in]{borders=aucune}{#1}}&
		\raisebox{-\depth}{\cemaththemeicon[height=1in]{borders=11}{#1}}&
		\raisebox{-\depth}{\cemaththemeicon[height=1in]{borders=10}{#1}}&
		\raisebox{-\depth}{\cemaththemeicon[height=1in]{borders=01}{#1}}&
		\raisebox{-\depth}{\cemaththemeicon[height=1in]{borders=00}{#1}}\xdef\tmpcptxint{\inteval{\tmpcptxint+5}}\\ \hline
	}
\end{longtable}




\end{document}